\section{Nonlinear-regime study and certificate conservatism}
\label{sec:transport-followup}

We test whether valid confidence transport produces useful certificates and
decisions in a controlled nonlinear model.  The prespecified follow-up evaluates
600 policy trajectories; a separate 294 development and tuning trajectories
precede evaluation.  The original near-linear study is retained separately in
\Cref{sec:transport-instantiation}; its evidence is not pooled with this study.

\paragraph{Model and prespecified protocol.}
We use Gaussian observations with mean
$\mu_{\btheta}(z)=\sqrt W\tanh(\phi(z)^\top\btheta/\sqrt W)$,
with $\|\phi(z)\|_2\le1$ and projected parameters $\|\btheta\|_2\le1$.
Contexts have dimension four, there are five actions, and the context--action
feature map has dimension 29.  Noise has standard deviation $0.25$, the metric
ridge is one, and the confidence parameter is $\delta=0.05$.
A fixed teacher generates the rewards.  We fix
$T=1000$ and $W\in\{16,1,0.25\}$, with 50 common evaluation seeds per condition
and four policies.  Changing $W$ changes both the reward function and its
curvature; cross-condition trends do not isolate a causal effect of metric
drift.  The conditions were specified before their evaluation outcomes were
observed.  Development, tuning, and evaluation use disjoint seed sets.

All policies use the corrected center and confidence radius of
\Cref{cor:corrected-center-new}.  \emph{Hessian} uses the analytic path
certificate $\bar D_t^Q$; \emph{endpoint} substitutes the dense endpoint Thompson
distance; \emph{frozen reference} uses the reference width directly; and
\emph{naive current} uses the current width without transport, an uncertified
control.  All four policies update their representations: ``frozen'' describes
the collection-time metric, not a fixed representation.  They share contexts
and potential-noise streams but have their own adaptive action histories.
The common optimizer, learning rate $3\times10^{-4}$ with 20 full-batch steps
followed by projection each round, was selected from nine candidates by prediction MSE
on uniform-random tuning histories.  This controls one implementation choice;
it does not establish policy-specific regret-optimal tuning.

Dense float64 matrices and Cholesky solves implement the exact-operator
abstraction, and finite action enumeration uses the smallest-index tie rule.
This is a 29-parameter diagnostic study, not a test of approximate solvers,
large-model scalability, or real-data performance.  Further protocol and
provenance details appear in \Cref{app:transport-followup}.

\paragraph{Prespecified policy outcomes.}
Hessian transport has higher mean cumulative pseudo-regret than every
comparator in every condition (\Cref{fig:transport-followup-outcomes});
\Cref{tab:transport-followup-outcomes} gives the numerical values.  All nine
paired 95\% intervals for comparator-minus-Hessian regret exclude zero.
These are pointwise intervals from 10,000 paired seed-bootstrap resamples per
comparison, without multiplicity adjustment.  They quantify evaluation-seed
variation conditional on the fixed teacher, design, and selected optimizer;
they do not include uncertainty from environment selection or tuning.
The comparisons are total-policy outcomes on matched exogenous streams, not
fixed-history interventions.  They establish neither a universal policy
ordering nor a causal attribution to the certificate alone.

For descriptive context, the online context-free mean-only benchmark has mean
regrets 185.26, 183.06, and 176.67 at $W=16,1,0.25$, respectively; the best fixed
action in hindsight has mean regrets 180.01, 177.97, and 171.84.  These means
use the same evaluation streams.  Hessian's sample means are close to the
context-free benchmark's, but no equivalence test was performed, and similar
regret does not imply matching actions or identical use of context.
\Cref{app:transport-followup} defines both benchmarks.

\begin{figure}[p]
  \centering
  % Generated by tools/make_transport_explanatory_figures.py; no experiments rerun.
% Original input SHA-256 f28acc2d601d7524f5b30a55f396c5408a7f37a26fd36bde49a3477024bc6658  review/transport_review_followup/full-cce-full-local/aggregate_compact.json
% Original input SHA-256 fb3293daacba3511ca5f3269d67e9aa0f0aca2708ad9876661ca062e197a59a4  review/transport_review_followup/full-cce-full-local/drift_and_reference_diagnostics.json
% Original input SHA-256 b1a6e9c5d8cfe70d3c3e3bb7cbd818502c21bb3b2f276426702cb016db323ecc  review/transport_review_followup/full-cce-full-local/terminal_bound_decomposition.json
\definecolor{exHessian}{HTML}{D55E00}
\definecolor{exEndpoint}{HTML}{0072B2}
\definecolor{exFrozen}{HTML}{009E73}
\definecolor{exNaive}{HTML}{CC79A7}
\begin{tikzpicture}
\begin{groupplot}[
  group style={group size=2 by 3,horizontal sep=0.42in,vertical sep=0.60in},
  scale only axis,width=2.30in,height=1.17in,
  /tikz/font=\normalsize,tick label style={font=\normalsize},
  label style={font=\normalsize},title style={font=\normalsize,align=center},
  axis lines=left,tick align=outside,clip=false,
  ymin=-0.5,ymax=3.5,ytick={0,1,2,3},
  yticklabels={Naive current,Frozen reference,Endpoint,Hessian},
  xmajorgrids,grid style={black!12},
]
\nextgroupplot[title={(a) Mean cumulative\\pseudo-regret},xmin=145,xmax=190,xtick={150,170,190}]
\node[anchor=south east] at (rel axis cs:1,1.04) {$W=16$};
\addplot[only marks,mark=*,mark size=2.8pt,color=exHessian] coordinates {(185.10060051413322,3)};
\addplot[only marks,mark=triangle*,mark size=2.8pt,color=exEndpoint] coordinates {(152.59341955672849,2)};
\addplot[only marks,mark=square*,mark size=2.8pt,color=exFrozen] coordinates {(151.78460762261437,1)};
\addplot[only marks,mark=diamond*,mark size=2.8pt,color=exNaive] coordinates {(152.74572489523797,0)};
\nextgroupplot[title={(b) Comparator minus Hessian\\Pointwise 95\% paired CI},xmin=-36,xmax=2,xtick={-30,-20,-10,0},yticklabels={,,,}]
\draw[densely dashed,black!65] (axis cs:0,-0.5) -- (axis cs:0,3.5);
\node[anchor=east] at (rel axis cs:0.90,0.88) {Hessian $=0$};
\addplot[no marks,line width=1pt,color=exEndpoint] coordinates {(-33.032610627919425,2) (-31.99704060364553,2)};
\addplot[no marks,color=exEndpoint,line width=1pt] coordinates {(-33.032610627919425,1.9099999999999999) (-33.032610627919425,2.0899999999999999)};
\addplot[no marks,color=exEndpoint,line width=1pt] coordinates {(-31.99704060364553,1.9099999999999999) (-31.99704060364553,2.0899999999999999)};
\addplot[only marks,mark=triangle*,mark size=1.5pt,color=exEndpoint] coordinates {(-32.507180957404692,2)};
\addplot[no marks,line width=1pt,color=exFrozen] coordinates {(-33.787557878165572,1) (-32.860633812570086,1)};
\addplot[no marks,color=exFrozen,line width=1pt] coordinates {(-33.787557878165572,0.91000000000000003) (-33.787557878165572,1.0900000000000001)};
\addplot[no marks,color=exFrozen,line width=1pt] coordinates {(-32.860633812570086,0.91000000000000003) (-32.860633812570086,1.0900000000000001)};
\addplot[only marks,mark=square*,mark size=1.5pt,color=exFrozen] coordinates {(-33.315992891518853,1)};
\addplot[no marks,line width=1pt,color=exNaive] coordinates {(-32.827710367754008,0) (-31.904417058504045,0)};
\addplot[no marks,color=exNaive,line width=1pt] coordinates {(-32.827710367754008,-0.089999999999999997) (-32.827710367754008,0.089999999999999997)};
\addplot[no marks,color=exNaive,line width=1pt] coordinates {(-31.904417058504045,-0.089999999999999997) (-31.904417058504045,0.089999999999999997)};
\addplot[only marks,mark=diamond*,mark size=1.5pt,color=exNaive] coordinates {(-32.354875618895228,0)};
\nextgroupplot[xmin=145,xmax=190,xtick={150,170,190}]
\node[anchor=south east] at (rel axis cs:1,1.04) {$W=1$};
\addplot[only marks,mark=*,mark size=2.8pt,color=exHessian] coordinates {(182.90041423907414,3)};
\addplot[only marks,mark=triangle*,mark size=2.8pt,color=exEndpoint] coordinates {(173.18693942935238,2)};
\addplot[only marks,mark=square*,mark size=2.8pt,color=exFrozen] coordinates {(173.19525926714039,1)};
\addplot[only marks,mark=diamond*,mark size=2.8pt,color=exNaive] coordinates {(173.05176063754888,0)};
\nextgroupplot[xmin=-36,xmax=2,xtick={-30,-20,-10,0},yticklabels={,,,}]
\draw[densely dashed,black!65] (axis cs:0,-0.5) -- (axis cs:0,3.5);
\node[anchor=east] at (rel axis cs:0.90,0.88) {Hessian $=0$};
\addplot[no marks,line width=1pt,color=exEndpoint] coordinates {(-10.019311312014434,2) (-9.3891319855430968,2)};
\addplot[no marks,color=exEndpoint,line width=1pt] coordinates {(-10.019311312014434,1.9099999999999999) (-10.019311312014434,2.0899999999999999)};
\addplot[no marks,color=exEndpoint,line width=1pt] coordinates {(-9.3891319855430968,1.9099999999999999) (-9.3891319855430968,2.0899999999999999)};
\addplot[only marks,mark=triangle*,mark size=1.5pt,color=exEndpoint] coordinates {(-9.7134748097217418,2)};
\addplot[no marks,line width=1pt,color=exFrozen] coordinates {(-9.9877028811556823,1) (-9.4168901482167477,1)};
\addplot[no marks,color=exFrozen,line width=1pt] coordinates {(-9.9877028811556823,0.91000000000000003) (-9.9877028811556823,1.0900000000000001)};
\addplot[no marks,color=exFrozen,line width=1pt] coordinates {(-9.4168901482167477,0.91000000000000003) (-9.4168901482167477,1.0900000000000001)};
\addplot[only marks,mark=square*,mark size=1.5pt,color=exFrozen] coordinates {(-9.7051549719337054,1)};
\addplot[no marks,line width=1pt,color=exNaive] coordinates {(-10.190414583928037,0) (-9.5037094084246316,0)};
\addplot[no marks,color=exNaive,line width=1pt] coordinates {(-10.190414583928037,-0.089999999999999997) (-10.190414583928037,0.089999999999999997)};
\addplot[no marks,color=exNaive,line width=1pt] coordinates {(-9.5037094084246316,-0.089999999999999997) (-9.5037094084246316,0.089999999999999997)};
\addplot[only marks,mark=diamond*,mark size=1.5pt,color=exNaive] coordinates {(-9.8486536015252462,0)};
\nextgroupplot[xlabel={Mean regret; no uncertainty bars},xmin=145,xmax=190,xtick={150,170,190}]
\node[anchor=south east] at (rel axis cs:1,1.04) {$W=0.25$};
\addplot[only marks,mark=*,mark size=2.8pt,color=exHessian] coordinates {(176.4391375795567,3)};
\addplot[only marks,mark=triangle*,mark size=2.8pt,color=exEndpoint] coordinates {(172.08884041017453,2)};
\addplot[only marks,mark=square*,mark size=2.8pt,color=exFrozen] coordinates {(171.87558013200885,1)};
\addplot[only marks,mark=diamond*,mark size=2.8pt,color=exNaive] coordinates {(172.58350759875503,0)};
\nextgroupplot[xlabel={Difference in mean regret},xmin=-36,xmax=2,xtick={-30,-20,-10,0},yticklabels={,,,}]
\draw[densely dashed,black!65] (axis cs:0,-0.5) -- (axis cs:0,3.5);
\node[anchor=east] at (rel axis cs:0.90,0.88) {Hessian $=0$};
\addplot[no marks,line width=1pt,color=exEndpoint] coordinates {(-5.10915683531731,2) (-3.5723861655103155,2)};
\addplot[no marks,color=exEndpoint,line width=1pt] coordinates {(-5.10915683531731,1.9099999999999999) (-5.10915683531731,2.0899999999999999)};
\addplot[no marks,color=exEndpoint,line width=1pt] coordinates {(-3.5723861655103155,1.9099999999999999) (-3.5723861655103155,2.0899999999999999)};
\addplot[only marks,mark=triangle*,mark size=1.5pt,color=exEndpoint] coordinates {(-4.3502971693821095,2)};
\addplot[no marks,line width=1pt,color=exFrozen] coordinates {(-5.6421526276077403,1) (-3.4859126782584213,1)};
\addplot[no marks,color=exFrozen,line width=1pt] coordinates {(-5.6421526276077403,0.91000000000000003) (-5.6421526276077403,1.0900000000000001)};
\addplot[no marks,color=exFrozen,line width=1pt] coordinates {(-3.4859126782584213,0.91000000000000003) (-3.4859126782584213,1.0900000000000001)};
\addplot[only marks,mark=square*,mark size=1.5pt,color=exFrozen] coordinates {(-4.5635574475478728,1)};
\addplot[no marks,line width=1pt,color=exNaive] coordinates {(-4.6937096979062423,0) (-3.0324461345366807,0)};
\addplot[no marks,color=exNaive,line width=1pt] coordinates {(-4.6937096979062423,-0.089999999999999997) (-4.6937096979062423,0.089999999999999997)};
\addplot[no marks,color=exNaive,line width=1pt] coordinates {(-3.0324461345366807,-0.089999999999999997) (-3.0324461345366807,0.089999999999999997)};
\addplot[only marks,mark=diamond*,mark size=1.5pt,color=exNaive] coordinates {(-3.8556299808016607,0)};
\end{groupplot}
\end{tikzpicture}
\par\medskip
\noindent\begin{minipage}{\linewidth}
\raggedright
$T=1000$; 50 paired seeds per condition; separate adaptive histories.
All nine intervals: 10,000 paired bootstrap resamples, pointwise 95\%,
without multiplicity adjustment. Negative differences favor the comparator.
The $W$ conditions change the reward task as well as curvature.
\end{minipage}

  \caption{Policy outcomes and paired comparisons in every prespecified
  condition.  Left: descriptive mean cumulative pseudo-regret at $T=1000$,
  without error bars.  Right: each comparator's regret minus Hessian's, with
  pointwise 95\% paired seed-bootstrap intervals from 10,000 resamples of
  50 seeds per condition.  Negative differences favor the comparator.
  The intervals condition on the fixed teacher, design and selected optimizer;
  they are not multiplicity-adjusted.  Each policy follows its own adaptive
  history on matched exogenous streams.  Changing $W$ changes the reward task
  as well as curvature, so the rows do not isolate an effect of drift.}
  \label{fig:transport-followup-outcomes}
\end{figure}

\paragraph{Realized drift, validity diagnostics, and vacuous bounds.}
The primary policy reaches endpoint distances above one in 16 of 50 seeds at
$W=0.25$, including 286 of 50,000 seed--round observations
(\Cref{tab:transport-followup-drift-bounds}).  These round counts describe
dependent observations, not additional independent replications.  Target
labels used to encode $W$ are not realized distances.
All 600 evaluation trajectories, including naive current, pass the recorded
simultaneous reference-confidence, method-optimism, and deterministic
diagnostics.  Each policy--condition cell has 50 successes out of 50, with an
exact 95\% Clopper--Pearson interval of approximately $[92.9\%,100\%]$.
This observation does not establish population coverage, validate the theorem,
or show that transport is necessary.  Float64 checks and converged quadrature
are numerical diagnostics, not verified enclosures.

The primary bounds are dominated by path inflation in all three conditions.
Every primary-policy bound exceeds the deterministic regret cap
$2T\min(1,\sqrt W)$, which is 2000, 2000, and 1000 for the three conditions.
The frozen policy's own reference-geometry bounds also exceed these caps in
all 150 of its trajectories.  Their mean right-hand sides are much smaller
than the primary bounds but remain vacuous.
\Cref{fig:transport-followup-bound-scale} makes their scale visible relative
to the cap; values above the cap provide no improvement over that deterministic
upper bound.  Removing the path term along
Hessian's histories would be an accounting decomposition, not the bound or
outcome of a frozen policy following different histories.  Base dominance
reported for the original near-linear grid does not describe this follow-up.

\begin{figure}[t]
 \centering
 % Generated by tools/make_transport_explanatory_figures.py; no experiments rerun.
% Original input SHA-256 f28acc2d601d7524f5b30a55f396c5408a7f37a26fd36bde49a3477024bc6658  review/transport_review_followup/full-cce-full-local/aggregate_compact.json
% Original input SHA-256 fb3293daacba3511ca5f3269d67e9aa0f0aca2708ad9876661ca062e197a59a4  review/transport_review_followup/full-cce-full-local/drift_and_reference_diagnostics.json
% Original input SHA-256 b1a6e9c5d8cfe70d3c3e3bb7cbd818502c21bb3b2f276426702cb016db323ecc  review/transport_review_followup/full-cce-full-local/terminal_bound_decomposition.json
\definecolor{exHessian}{HTML}{D55E00}
\definecolor{exEndpoint}{HTML}{0072B2}
\definecolor{exFrozen}{HTML}{009E73}
\definecolor{exNaive}{HTML}{CC79A7}
\begin{tikzpicture}
\begin{axis}[
  scale only axis,width=5.55in,height=2.40in,
  /tikz/font=\normalsize,tick label style={font=\normalsize},
  label style={font=\normalsize},
  axis lines=left,tick align=outside,clip=false,
  xmin=-2,xmax=85,xtick={0,20,40,60,80},
  ymin=-0.4,ymax=2.6,ytick={0,1,2},yticklabels={$W=0.25$,$W=1$,$W=16$},
  xlabel={$\log_{10}(\text{arithmetic mean terminal RHS}/\text{cap})$},
  xmajorgrids,grid style={black!12},
  legend style={draw=none,at={(0.5,1.17)},anchor=south,legend columns=2,
    font=\normalsize,/tikz/every even column/.append style={column sep=12pt}},
]
\draw[densely dashed,black!65] (axis cs:0,-0.4) -- (axis cs:0,2.5);
\node[anchor=south west] at (axis cs:0,2.5) {cap};
\addplot[only marks,mark=*,mark size=3.2pt,color=exHessian] coordinates {(9.5127767175073252,2.1099999999999999) (37.186731520124027,1.1100000000000001) (80.909706296459532,0.11)};
\addlegendentry{Hessian}
\addplot[only marks,mark=square*,mark size=3.2pt,color=exFrozen] coordinates {(1.2746038084791189,1.8899999999999999) (1.7801161507515202,0.89000000000000001) (2.3633154176506039,-0.11)};
\addlegendentry{Frozen reference}
\end{axis}
\end{tikzpicture}
\par\medskip
\noindent\begin{minipage}{\linewidth}
\raggedright
$T=1000$; arithmetic means over each policy's own 50 histories.
The axis shows the logarithm of a mean ratio, not the mean of logarithms.
Positive values exceed the deterministic cap. Caps are 2000, 2000 and 1000
at $W=16,1,0.25$, respectively. No uncertainty intervals are plotted.
\end{minipage}

 \caption{The reported bound scale exceeds the deterministic regret cap.
 Each point is the base-ten logarithm of the mean terminal theorem RHS divided
 by $2T\min(1,\sqrt W)$, using the named policy's own 50 histories.
 The zero line represents equality with the cap; positive values indicate
 that the mean bound is larger.  These are descriptive means without
 uncertainty intervals.  Every individual bound also exceeds its cap, as
 checked in the retained diagnostics.  The comparison is not a decomposition
 of one policy's bound on another policy's history.}
 \label{fig:transport-followup-bound-scale}
\end{figure}

\begin{table}[t]
  \centering
  \caption{Prespecified drift and bound diagnostics.  The maximum and threshold
  counts use Hessian's 50,000 seed--round observations per condition; thresholds
  use $\ge$.  Mean terminal right-hand sides (RHS) use each named policy's own
  50 histories.  All individual primary and frozen-policy bounds exceed the
  corresponding deterministic regret cap.}
  \label{tab:transport-followup-drift-bounds}
  % Generated by tools/make_transport_followup_paper.py; no experiments rerun.
% Input SHA-256 914013bf6436e82a50865d5ce3570b4195e7cdc2e4f25b436eb5b28a9b99ece3  review/post_full_fable/independent/analysis.json
% Input SHA-256 fd67cdc031c02876ed1e18ea0444a35348b4cb80bdd261b3fd3d75e6e920b29e  review/transport_relative_diagnostics/full/analysis.json
% Input SHA-256 9d7939364bfd7737df5daaf80d43c0673828505466ede181fe31f42fe2f98bea  review/transport_relative_diagnostics/full/per_seed.jsonl
% Input SHA-256 f28acc2d601d7524f5b30a55f396c5408a7f37a26fd36bde49a3477024bc6658  review/transport_review_followup/full-cce-full-local/aggregate_compact.json
% Input SHA-256 fb3293daacba3511ca5f3269d67e9aa0f0aca2708ad9876661ca062e197a59a4  review/transport_review_followup/full-cce-full-local/drift_and_reference_diagnostics.json
% Input SHA-256 9b6bda78c88074e46e1ff3c4253b9623bb95f467f068b6339bee9d83e4e4fce1  review/transport_review_followup/full-cce-full-local/round_drift_summaries.jsonl
% Input SHA-256 e1c4935f879df7418a2a115692113d90e50d6e50e3819cbe3ad29ce40a9afc60  review/transport_review_followup/full-cce-full-local/seed_summaries.jsonl
% Input SHA-256 b1a6e9c5d8cfe70d3c3e3bb7cbd818502c21bb3b2f276426702cb016db323ecc  review/transport_review_followup/full-cce-full-local/terminal_bound_decomposition.json
% Input SHA-256 e8e0c87f16e2bf9252b27d022a473a114c33d395cc20303d156013568e9dc836  tools/make_transport_followup_paper.py
\begin{tabular}{@{}rrrrrr@{}}
\toprule
$W$ & \shortstack{Max.\\$d_{\rm Th}$} & \shortstack{Rounds at least\\$0.01\,/\,0.1\,/\,1$} & \shortstack{Mean RHS\\Hessian} & \shortstack{Mean RHS\\frozen policy} & Cap \\
\midrule
16 & 0.045 & 1495/0/0 & $6.51\times10^{12}$ & $37\,639$ & 2000 \\
1 & 0.627 & 49789/3804/0 & $3.07\times10^{40}$ & $120\,544$ & 2000 \\
0.25 & 1.581 & 49863/45623/286 & $8.12\times10^{83}$ & $230\,842$ & 1000 \\
\bottomrule
\end{tabular}

\end{table}

\paragraph{Exploratory decision analysis.}
After observing these outcomes, we independently recomputed score diagnostics
from all 600 preserved trajectories.  On Hessian's own histories, its selected
action equals the current-width argmax on 99.382\%, 99.962\%, and 99.982\% of
rounds for $W=16,1,0.25$.  A stronger sufficient condition---the bonus gap
between the two largest widths exceeding the entire range of prediction
centers---holds on 97.104\%, 99.372\%, and 99.508\% of rounds.
These pooled fractions describe extensive width-dominated selection in the
visited states, not an adaptive causal mechanism or an independent-round test.
The per-seed mechanism table in \Cref{app:transport-followup} applies this same
Hessian score to all four behavior histories; it does not report width dominance
of each behavior policy's own score.  The appendix also compares alternative
score recommendations on fixed histories and checks early numerical ranking
sensitivity.  The fixed-history comparisons respect
\Cref{prop:reference-width-floor}: tightening a transport certificate does not
create additional statistical information or order policies' adaptive regret.

\paragraph{Exploratory tighter endpoint certificates.}
The relative-weight construction in \Cref{cor:tanh-relative-weights} permits
uniform and category-specific endpoint comparisons without using the Hessian
path envelope.  We audit these comparisons on the same 600 recorded histories,
replaying representation updates with the recorded actions and rewards.
On primary histories, the category-specific outward factor gives mean
effective-width overheads of $0.28\%$--$12.68\%$ relative to the direct reference
width, versus $0.16\%$--$12.13\%$ for the dense endpoint factor across conditions
(\Cref{app:transport-followup}).  These are means of actionwise ratios, not
ratios of mean widths.  This analysis evaluates new scores
on existing states; it does not evaluate a new adaptive policy or demonstrate
an end-to-end computational advantage.  The relative certificate bounds an
endpoint comparison, not the original interpolated factor-path length, and
does not remove the historical and current approximation-error charges.
